\documentclass[12pt]{article}
\usepackage[utf8]{inputenc}
\usepackage[T1]{fontenc}
\usepackage{hyperref}
\title{Seimat Language Resources for K-12 Curriculum Development}
\author{Compiled by Claude Sonnet 5.5 (claude-sonnet-5-5) for Andr\'as Kornai}
\date{}
\begin{document}
\maketitle

\section{Introduction}
Seimat (ISO 639-3: ssg; Glottocode seim1238) is an Austronesian language of the Ninigo and Anchorite island groups in Manus Province, Papua New Guinea. It belongs to the Western Admiralty Islands group, together with Wuvulu-Aua and the extinct Kaniet. Estimates are about 1,000 speakers (Wikipedia, citing 1992) to about 1,200 (Omniglot). Searches found little published material. No dictionary, primer or textbook turned up.

\section{Primary Educational Resources}
No primers or textbooks were found in search results. Omniglot states that Seimat is taught in some schools and that a Latin-based writing system was proposed at a 2003 meeting of Seimat speakers from Pihon Island. It adds that other spelling systems exist. I did not find a school programme or curriculum documents.

\section{Linguistic and Reference Materials}
\begin{itemize}
\item Beata Wozna and Teresa Wilson, \emph{Seimat Grammar Essentials} (2005), SIL Data Papers on Papua New Guinea Languages, vol.\ 48, a grammar sketch of about 96 pages. Only the catalogue record was found, not a digital copy: \url{https://glottolog.org/resource/reference/id/42399}
\item Seimat numeral system data (Channumerals database): \url{https://lingweb.eva.mpg.de/channumerals/Seimat.htm}
\item Omniglot page on Seimat, covering writing and orthography: \url{https://www.omniglot.com/writing/seimat.htm}
\item Wikipedia overview of Seimat: \url{https://en.wikipedia.org/wiki/Seimat_language}
\end{itemize}

\section{Religious and Community Resources}
No scripture translations or community organisations were found. The Joshua Project profile for the Ninigo and Seimat people gives demographic context only: \url{https://joshuaproject.net/people_groups/13405/PP}

\section{Digital and Audio Resources}
\begin{itemize}
\item Verbix Seimat verb list (a short word list, not a dictionary): \url{https://docs.verbix.com/Languages/Seimat}
\item Rosetta Panglossia entry for ssg: \url{https://rosettapanglossia.longnow.org/index.php?title=Ssg}
\end{itemize}
The search summaries mention that Kaipuleohone holds recordings and written materials for Seimat, but no URL for that archive was found, so none is cited here.

\section{Curriculum Development Context}
Little published material exists. The Seimat grammar sketch and the orthography notes are the main foundations. A vernacular curriculum would need new work with Seimat speakers, such as vocabulary, stories and an agreed spelling, and should be checked with SIL Papua New Guinea and Manus Province education authorities. This search did not reach those bodies.

\end{document}